\subsection{Stellar algebra of a locally compact group}

DEFINITION 9. --- Let G be a locally compact group and let A be the involutive Banach algebra of bounded measures on G admitting a density with respect to a Haar measure on G (Example 4 of I, p.~99). The enveloping stellar algebra of the involutive Banach algebra A is called the stellar algebra of G. It is denoted by Stell(G).

Note. --- Let \(\nu\) be a left Haar measure on G and let \(\Delta\) be its modulus. The map \(f \mapsto f \cdot \nu\) is an isometric isomorphism of the algebra \(\mathrm{L}^{1}(\mathrm{G}, \nu)\) onto A (loc. cit.). One can therefore also define Stell(G) as the enveloping stellar algebra of the normed involutive algebra \(\mathrm{L}^{1}(\mathrm{G}, \nu)\) .

Let G be a locally compact group and let \(\nu\) be a left Haar measure on G. For \(\mu \in \mathcal{M}^1 (\mathrm{G})\) and \(f\in \mathrm{L}^2 (\mathrm{G},\nu)\) , one has \(\mu *f\in \mathrm{L}^2 (\mathrm{G},\nu)\) (INT, VIII, §4, prop. 6). Let us then denote by \(\pmb{\gamma}(\mu)\) the endomorphism \(f\mapsto \mu *f\) of \(\mathrm{L}^2 (\mathrm{G},\nu)\) . The map \(\mu \mapsto \pmb{\gamma}(\mu)\) is a representation of the algebra \(\mathcal{M}^{1}(\mathrm{G})\) into the Banach algebra \(\mathcal{L}(\mathrm{L}^{2}(\mathrm{G},\nu))\) of continuous endomorphisms of \(\mathrm{L}^2 (\mathrm{G},\nu)\) (INT, VIII, §4, cor. of prop. 6). On the other hand, \(\gamma (\breve{\mu})\) is the transpose of the endomorphism \(\gamma (\mu)\) (INT, VIII, §4, no. 3, prop. 8). It follows that \(\gamma (\mu^{*})\) is the adjoint endomorphism of \(\gamma (\mu)\) , and therefore that the map \(\gamma \colon \mu \mapsto \gamma (\mu)\) is a morphism of involutive algebras from \(\mathcal{M}^{1}(\mathrm{G})\) into the stellar algebra \(\mathcal{L}(\mathrm{L}^{2}(\mathrm{G},\nu))\) , called the left regular representation of \(\mathcal{M}^{1}(\mathrm{G})\) in \(\mathrm{L}^2 (\mathrm{G},\nu)\) . According to INT, VIII, §4, no. 7, prop. 19, this representation is faithful.

Let \(j\) be the canonical map from \(\mathrm{L}^{1}(\mathrm{G},\nu)\) into \(\mathrm{Stell}(\mathrm{G})\) . By restriction to \(\mathrm{L}^{1}(\mathrm{G},\nu)\) , the regular representation \(\pmb{\gamma}\) defines an injective morphism of involutive algebras from \(\mathrm{L}^{1}(\mathrm{G},\nu)\) into \(\mathcal{L}(\mathrm{L}^{2}(\mathrm{G},\nu))\) , called the left regular representation of \(\mathrm{L}^{1}(\mathrm{G})\) in \(\mathrm{L}^{2}(\mathrm{G})\) . According to prop. 20, there exists a unique morphism \(\pmb{\gamma}'\colon \mathrm{Stell}(\mathrm{G})\to \mathcal{L}(\mathrm{L}^{2}(\mathrm{G},\nu))\) such that \(\pmb{\gamma} = \pmb{\gamma}'\circ j\) . We say that \(\pmb{\gamma}'\) is the left regular representation of \(\mathrm{Stell}(\mathrm{G})\) in \(\mathrm{L}^{2}(\mathrm{G},\nu)\) . By abuse of notation, we will still denote it by

\[
\boldsymbol {\gamma} ^ {\prime} (\varphi) (f) = \varphi * f\tag{7}
\]

for \(f\in \mathrm{L}^2 (\mathrm{G},\nu)\) and \(\varphi \in \mathrm{Stell}(\mathrm{G})\) . We have

\[
\| \varphi * f \| _ {2} \leqslant \| \varphi \| _ {*} \| f \| _ {2}.\tag{8}
\]

Remark. --- In general, the left regular representation \(\gamma'\) of Stell(G) in \(L^{2}(G,\nu)\) is not faithful. One can show that it is faithful if and only if there exists a positive linear form f on \(\mathrm{L}_{\mathbf{R}}^{\infty}(\mathrm{G})\) such that \(f(1)=1\) and \(f(\boldsymbol{\gamma}(g)x)=f(x)\) for all \((g,x)\in\mathrm{G}\times\mathrm{G}\) (we then say that the group G is amenable, cf.~EVT, IV, p.~73, exercise 4).

PROPOSITION 22. --- The canonical map j from \(\mathrm{L}^{1}(\mathrm{G},\nu)\) into Stell(G) is injective and has dense image.

The image of \(j\) is dense by definition of the enveloping stellar algebra of a normed involutive algebra. Since the left regular representation \(\gamma\) is faithful, the injectivity of \(j\) follows from the equality \(\gamma = \gamma' \circ j\) .

COROLLARY. --- The algebra \(\mathrm{L}^{1}(\mathrm{G},\nu)\) is radical-free.

This follows from prop. 21 of I, p.~125 and from prop. 22.

Thus one can identify \(\mathrm{L}^{1}(\mathrm{G},\nu)\) with a dense involutive subalgebra of \(\mathrm{Stell}(\mathrm{G})\) , and the canonical injection of \(\mathrm{L}^{1}(\mathrm{G},\nu)\) into \(\mathrm{Stell}(\mathrm{G})\) is then continuous.

Suppose now that the group G is unimodular (INT, VII, §1, no. 3, def. 3). One can then repeat the same arguments starting from the right regular representation \((f,\mu)\mapsto\boldsymbol{\delta}(\mu)(f)=f*\breve{\mu}\) of \(\mathrm{L}^{2}(\mathrm{G},\nu)\times\mathcal{M}^{1}(\mathrm{G})\) into \(\mathrm{L}^{2}(\mathrm{G},\nu)\) . We then define a morphism \(\delta^{\prime}\) of \(\operatorname{Stell}(\mathrm{G})\) into \(\mathcal{L}(\mathrm{L}^{2}(\mathrm{G},\nu))\) such that \(\delta=\delta^{\prime}\circ j\) , and define \(\boldsymbol{\delta}^{\prime}(\varphi)(f)=f*\varphi\) for \(f\in\mathrm{L}^{2}(\mathrm{G},\nu)\) and \(\varphi\in\operatorname{Stell}(\mathrm{G})\) .

For \(\varphi, \psi \in \text{Stell}(G)\) , one has \(\boldsymbol{\delta}'(\psi) \circ \boldsymbol{\gamma}'(\varphi) = \boldsymbol{\gamma}'(\psi) \circ \boldsymbol{\delta}'(\varphi)\) , that is,

\[
(\varphi * f) * \psi = \varphi * (f * \psi)\tag{9}
\]

for all \(f \in \mathrm{L}^{2}(\mathrm{G}, \nu)\) . Indeed, this formula is true for \(\varphi, \psi \in \mathrm{L}^{1}(\mathrm{G}, \nu)\) , and the maps \((\varphi, \psi) \mapsto (\varphi * f) * \psi\) and \((\varphi, \psi) \mapsto \varphi * (f * \psi)\) are continuous bilinear maps from \(\operatorname{Stell}(\mathrm{G}) \times \operatorname{Stell}(\mathrm{G})\) into \(\mathrm{L}^{2}(\mathrm{G}, \nu)\) .

